\subsection{De rang van een matrix}

\begin{stelling}[5.1.2]
    Zij $A \in M_{m,n}(K)$ met kolommen $\{A_1, \dots, A_n\}$ en rijen $\{R_1, \dots, R_m\}$. Dan is $\dim \operatorname{span}(A_1, \dots, A_n) = \dim \operatorname{span}(R_1, \dots, R_m)$.
\end{stelling}
\begin{proof}
    \(k = \dim(\operatorname{span}(A_1,\dots,A_n))\) we kunnen hiervoor een basis \(\mathcal{B} = \{b_1,\dots,b_k\}\) vinden met \(k\) elementen.

    \(A_j = \sum_i^k x_{ij} b_i \implies A = B \cdot X\). \(B \in M_{m,k}(K)\) en \(X \in M_{k,n}(K)\).

    \(R_i = \sum_j^m b_{ij} x_j \implies \dim(\operatorname{span}(R_1,\dots,R_n)) \leq \dim(\operatorname{span}(X_1,\dots,X_n))\).

    We weten ook dat \(\dim(\operatorname{span}(R_1,\dots,R_n)) \leq k = \dim(\operatorname{span}(A_1,\dots,A_n))\).

    Herhalen we dit voor de getransponeerde matrix,
    dan vinden we  \(\dim(\operatorname{span}(A_1,\dots,A_n)) \leq k = \dim(\operatorname{span}(R_1,\dots,R_n))\).

    Bijgevolg is \(\dim(\operatorname{span}(A_1,\dots,A_n)) = \dim(\operatorname{span}(R_1,\dots,R_n)) = k\)
\end{proof}

\begin{lemma}[5.1.6]
    Zij $A \in M_{m,n}(K)$ met $\operatorname{rk}(A)$ de rang van $A$.
    \begin{enumerate}
        \item[(i)] Beschouw de lineaire afbeelding $L_A \colon K^n \to K^m \colon v \mapsto Av$. Dan is
        \[ \operatorname{im}(L_A) = \operatorname{span}\{A_1, \dots, A_n\}; \]
        bijgevolg is $\operatorname{dim}(\operatorname{im} L_A) = \operatorname{rk}(A)$.
        \item[(ii)] Elementaire rijoperaties laten de rang van een matrix invariant.
        \item[(iii)] Zij $A \in M_{m,n}(K)$ en $B \in M_{n,r}(K)$. Dan is $\operatorname{rk}(AB) \le \operatorname{rk}(A)$ en $\operatorname{rk}(AB) \le \operatorname{rk}(B)$.
        \item[(iv)] Zij $A \in \operatorname{GL}_n(K)$ en $B \in \operatorname{Mat}_{n,r}(K)$. Dan is $\operatorname{rk}(AB) = \operatorname{rk}(B)$.
        \item[(v)] Zij $A \in \operatorname{GL}_n(K)$ en $B \in \operatorname{Mat}_{r,n}(K)$. Dan is $\operatorname{rk}(BA) = \operatorname{rk}(B)$.
    \end{enumerate}
\end{lemma}
\begin{proof}(i)
    Kies als basis de eenheidsbasis in \(K^n\). \(\operatorname{im}(L_A)\) wordt voortgeebracht door \(\{Ae_1,\dots,Ae_n\}\).
    \(Ae_i \in K^m\) is gelijk aan de \(i\)-de kolom van \(A\).
\end{proof}
\begin{proof}(ii)
      Elementaire rijoperaties laten de oplossingsverzameling invariant.
\end{proof}
\begin{proof}(iii)
    \(AB(v) = A(Bv) \in \operatorname{im}(A) \implies \operatorname{rk}(AB) \leq \operatorname{rk}(A) \)

    \(\operatorname{rk}(AB) \leq \operatorname{rk}(B) \) volgt door hetzelfde te doen met de getransponeerde.
\end{proof}
\begin{proof}(iv)
      \(B = A^{-1}AB \implies \operatorname{rk}(B) = \operatorname{rk}(A^{-1}AB) \leq \operatorname{rk}(AB)\).\\
      Met de ongelijkheid uit (iii) volgt nu \(\operatorname{rk}(B) = \operatorname{rk}(AB)\).
\end{proof}
\begin{proof}(v)
      Anagloog aan (iv).
\end{proof}

\subsection{Affiene deelruimten}

\begin{lemma}[5.2.4]
    Zij $V$ een vectorruimte, $v_1, v_2 \in V$ en $W_1, W_2$ deelruimten van $V$. Dan is $v_1 + W_1 = v_2 + W_2$ als en slechts als $W_1 = W_2$ en $v_1 - v_2 \in W_1$.
\end{lemma}
\begin{proof}
    \([\Rightarrow]\) Als \(v_1 + W_1 = v_2 + W_2\): Stel \(w_0 = v_1 -v_2\).

    Dan is \(W_1 + v_1-v_2= W_1 + v_0 = W_2\). Omdat \(0 \in W_2\), moet \(w_0 \in W_1\), maar dan \(W_1 + w_0 = W_1\).
    Bijgevolg is \(W_1 = W_2\).

    \([\Leftarrow]\) Stel omgekeerd $W_1 = W_2$ en $v_1 - v_2 \in W_1$:

    \(v_1 + W_1 = \{v_1 + w | w\in W_1\}\) en \(v_2 + W_2 = \{v_1 - w_0 + w | w\in W_2\}\)

    Het is duidelijk dat de verzamelingen \(\forall w \in W_1 \text{ en } \forall (w-w_0) \in W_1\) gelijk zijn.

    \(\implies v_1 + W_1 = v_2 + W_2\).
\end{proof}

\subsection{De oplossingsverzameling van een stelsel}

\begin{lemma}[5.3.1]
    Beschouw het stelsel $AX = w$ met $A \in M_{m,n}(K)$, $w \in K^m$ en $X$ een kolomvector met $n$ onbekenden. Beschouw de lineaire afbeelding
    \[ L_A \colon K^n \to K^m \colon v \mapsto Av. \]
    \begin{enumerate}
        \item[(i)] Het stelsel $AX = w$ is strijdig als en slechts als $w \notin \operatorname{im} L_A$.
        \item[(ii)] Als het stelsel $AX = w$ niet strijdig is, dan is de oplossingsverzameling een affiene deelruimte. Als $v_0 \in K^n$ een oplossing is, dan is de oplossingsverzameling gegeven door $L_A^{-1}(w) = v_0 + \operatorname{ker} L_A$.
        \item[(iii)] Als het stelsel homogeen is, m.a.w. van de vorm $AX = 0$, dan is de oplossingsverzameling gelijk aan de deelruimte $\operatorname{ker} L_A$.
    \end{enumerate}
\end{lemma}
De oplossingsverzameling \(AX = w\) is:
\[
\{v\in K^n | Av = w\} = \{v \in K^n | L_A(v) = w\} = L_A^{-1}(w).
\]
\begin{proof}(i)
    \(L_A^{-1}(w) = \emptyset \iff w \notin \operatorname{im} L_A\).
\end{proof}
\begin{proof}(ii)
    Als \(v_0\) een oplossing is, is \(v_0 + \ker L_A \) ook een oplossing, want \(L_A (v_0 + \ker L_A) = L_A(v_0) =  w\).
    
    \(\implies    v_0 + \ker L_A \subseteq L_A^{-1}(w) \).
    Om de gelijkheid aan te tonen moeten we tonen dat elke willekeurige oplossing in de affiene deelruimte ligt.

    Neem \(v \in K^n\):
    \(L_A (v_0 + v) = L_A (v_0)  + L_A ( v) = w + L_A (v)\). Waar we de lineariteit van \(L_A\) hebben toegepast.

    Deze oplossing voldoet pas als \(L_A (v) = 0 \implies v \in \ker L_A \).

    We concluderen dat \(L_A^{-1}(w) = v_0 + \operatorname{ker} L_A\).

\end{proof}
\begin{proof}(iii)
    De oplossingsverzameling is \(\{v \in K^n | L_A(v) = 0\} = L_A^{-1}(0)\).

    Per definitie is \(  L_A^{-1}(0) = \ker L_A\).
\end{proof}

\begin{stelling}[5.3.2]
    Beschouw het stelsel $AX = w$ met $A \in M_{m,n}(K)$, $w \in K^m$ en $X$ een kolomvector met $n$ onbekenden.
    \begin{enumerate}
        \item[(i)] Het stelsel heeft een oplossing als en slechts als $\operatorname{rk}(A) = \operatorname{rk}(A|w)$.
        \item[(ii)] Stel dat het stelsel niet strijdig is. Dan is de dimensie van de oplossingsverzameling gelijk aan $n - \operatorname{rk}(A)$.
    \end{enumerate}
\end{stelling}
\begin{proof}(i)
    Uit lemma 5.3.1 (i) heeft het stelsel \(AX = w\) (een) oplossing(en) \(\iff w \in \operatorname{im}L_A\).
    
    \(\operatorname{rk}(A|w) 
    = \dim \operatorname{im}(A|w)
    = \dim \operatorname{span}(A_1,\dots,A_n,w)
    = \dim \operatorname{span}(A_1,\dots,A_n)
    = \dim \operatorname{im}(A)
    = \operatorname{rk}(A)
     \)

     Want \(w = \sum_{i}^{n}\lambda_i A_i\).
\end{proof}
\begin{proof}(ii)
    De oplossingsverzameling van een affiene deelruimte is \(v_0 + \ker L_A\).

    \(\dim(v_0 + \ker L_A) = \dim (\ker L_A) = \dim (K^n) - \dim(\operatorname{im}L_A) = n - \operatorname{rk}(A)\).
\end{proof}

\begin{gevolg}[5.3.3]
    Beschouw het stelsel $AX = w$, $A \in M_n(K)$. Als $\operatorname{rk}(A) = n$, dan heeft het stelsel een unieke oplossing.
\end{gevolg}
\begin{proof}
    \(\dim(v_0 + \ker L_A) = n - \operatorname{rk}(A) = 0\).

    De dimensie van de oplossingsverzameling is \(0\). Dit wijst op een unieke oplossing.
\end{proof}

\subsection{ Determinanten}
\textbf{Hiervan zijn de bewijzen niet te kennen!}
\subsection{Inverteerbaarheid van matrices}
\textbf{Hiervan zijn de bewijzen niet te kennen! De stellingen, eigenschappen en rekenregels zijn wel belangrijk!}